\documentclass[../HDF5_RFC.tex]{subfiles}

\begin{document}

%% ======================================================================
%%  Section 7 – Built-in Filter Parameters
%% ======================================================================
\section{Built-in Filter Parameters}
\label{sec:builtin-filters}

Each built-in filter MUST implement \texttt{set\_config} and \texttt{get\_config} so that the
string API is fully functional without third-party plugins.

For built-in filters that have dedicated \texttt{H5Pset\_*} functions with non-trivial
validation logic (particularly SZIP and Scale-Offset), the \texttt{set\_config}
callback MUST delegate to the existing internal validation code rather than
reimplementing it independently. This prevents divergence between the two
configuration paths and ensures that validation constraints are applied
consistently regardless of whether the filter was configured via
\texttt{H5Pset\_szip} or \texttt{H5Pset\_filter\_by\_name}.

\verysubsection{Implementation note}
Built-in filters are currently registered internally as \texttt{H5Z\_class2\_t}
structs. Each built-in filter's internal registration MUST be converted to
\texttt{H5Z\_class3\_t} with \texttt{set\_config} and \texttt{get\_config} added directly to the
struct definition (e.g., in \texttt{H5Zdeflate.c}, \texttt{H5Zszip.c}, etc.).

This conversion is entirely internal to the library. External code never
directly accesses these struct definitions, so there is no public API
compatibility concern. The converted structs are passed to the existing
\texttt{H5Zregister()} path, which already handles \texttt{H5Z\_class3\_t} registrations
(\SectionRef{sec:class3}). No parallel mechanism or additional internal API is
required.

\subsection{Deflate}

\begin{itemize}
\item \texttt{level} (integer 0--9, default 6): Compression level.
\item Example: \texttt{"level=6"}
\end{itemize}

\subsection{Shuffle}

No parameters. \texttt{params} should be \texttt{NULL}.

\subsection{Fletcher32}

No parameters. \texttt{params} should be \texttt{NULL}.

\subsection{SZIP}

\begin{itemize}
\item \texttt{coding} (\texttt{"entropy"} or \texttt{"nn"}, default \texttt{"nn"}): Coding
  method.
\item \texttt{pixels\_per\_block} (even integer 2--32, default 16): Pixels per block.
\item The \texttt{set\_config} callback translates these into the appropriate option
  masks (\texttt{H5\_SZIP\_EC\_OPTION\_MASK} or \texttt{H5\_SZIP\_NN\_OPTION\_MASK}) and
  delegates to the existing \texttt{H5Pset\_szip} validation path. The existing
  \texttt{H5Pset\_szip} (\texttt{H5Pdcpl.c}) performs the following operations beyond
  simple parameter translation: checks whether the SZIP encoder is available
  (\texttt{H5Z\_FILTER\_CONFIG\_ENCODE\_ENABLED}); validates \texttt{pixels\_per\_block} is
  even and does not exceed \texttt{H5\_SZIP\_MAX\_PIXELS\_PER\_BLOCK}; applies
  \texttt{H5\_SZIP\_ALLOW\_K13\_OPTION\_MASK}; sets \texttt{H5\_SZIP\_RAW\_OPTION\_MASK};
  strips \texttt{H5\_SZIP\_CHIP\_OPTION\_MASK}, \texttt{H5\_SZIP\_LSB\_OPTION\_MASK}, and
  \texttt{H5\_SZIP\_MSB\_OPTION\_MASK}; and forces \texttt{H5Z\_FLAG\_OPTIONAL} regardless
  of the caller's flag choice. The \texttt{set\_config} callback MUST replicate all
  of this behavior. Consequently, the \texttt{flags} argument to
  \texttt{H5Pset\_filter\_by\_name} is overridden for SZIP---the filter is always
  optional. This is existing behavior of \texttt{H5Pset\_szip} and SHOULD be
  documented but is not changed by this feature. Callers that need to
  detect the flag override can compare the flags they passed with the
  flags returned by \texttt{H5Pget\_filter2} after the call.
\item Example: \texttt{"coding=entropy, pixels\_per\_block=8"}
\end{itemize}

\subsection{N-Bit}

No user-facing parameters (derived from datatype). \texttt{params} should be
\texttt{NULL}.

\subsection{Scale-Offset}

\begin{itemize}
\item \texttt{scale\_type} (\texttt{"float\_dscale"}, \texttt{"float\_escale"}, or
  \texttt{"int"}): Scale type. Required.
\item \texttt{scale\_factor} (integer): Scale factor. Required.
\item The \texttt{set\_config} callback delegates to the existing
  \texttt{H5Pset\_scaleoffset} validation path.
\item Example: \texttt{"scale\_type=float\_dscale, scale\_factor=4"}
\end{itemize}


%% ======================================================================
%%  Section 8 – Thread Safety and Concurrency
%% ======================================================================
\section{Thread Safety and Concurrency}
\label{sec:thread-safety}

\subsection{Global Registry Protection}

The canonical name registry is a global mutable data structure. In
thread-safe builds of HDF5 (configured with \texttt{--enable-threadsafe}), all
access to the registry---including lookups during \texttt{H5Pset\_filter\_by\_name}
and mutations during \texttt{H5Zregister\_name} or plugin auto-registration---MUST
be serialized.

The existing HDF5 global API lock serializes all public API entry points.
This is sufficient to prevent concurrent mutation of the registry from
user-facing calls. Plugin auto-registration during on-demand search
(\SectionRef{sec:name-discovery}) occurs within the same API call scope and is
therefore covered by the same lock. No new locking mechanisms are introduced.

\verysubsection{Future-proofing note}
\texttt{H5Zlookup\_by\_name} is a pure read operation with
no I/O side effects. If the library moves toward finer-grained concurrency
(as discussed in the context of VOLs and asynchronous I/O), the name registry
SHOULD be protected by a dedicated readers-writer (RW) lock rather than the
global API lock, allowing concurrent lookups without serialization. The
internal data structure chosen for the registry should be compatible with this
future migration.

In addition to the registry lookup, the on-demand plugin scan triggered by
\texttt{H5PL\_load} performs filesystem I/O (directory iteration and shared-library
opens). In thread-safe builds, this I/O is also serialized under the global
API lock. When many threads simultaneously encounter an unknown filter name at
startup, all threads queue behind the first thread performing the scan,
producing a convoy effect. Reducing lock granularity for the scan itself---for
example by allowing concurrent read-only scans while only serializing the
final registry insertion---would therefore benefit throughput independently of
the registry RW-lock migration described above.

\subsection{Parallel HDF5 (MPI) Considerations}

In Parallel HDF5, all MPI ranks participating in collective dataset creation
MUST present identical dataset creation property lists (DCPLs). Because filter
name resolution depends on locally loaded plugins, non-deterministic plugin
loading across ranks (e.g., due to filesystem latency or differing
\texttt{HDF5\_PLUGIN\_PATH} settings) could cause Rank~0 to resolve \texttt{"zfp"}
successfully while Rank~1 fails.

This is not a new problem introduced by this design---the same class of
issue exists today when any rank fails to load a required plugin, and there
are already several runtime-specific mechanisms (environment variables,
plugin path configuration, dynamic library search order) that can cause
per-rank divergence independently of whether filters are identified by name
or by ID.

An additional parallel-specific concern is that parameter string parsing is
performed independently on every rank, even though the resulting
\texttt{cd\_values} are identical across ranks and the work need only be done
once.  This is an inherent inefficiency of the current design, though it
should be noted that the existing plugin interface was not designed with
parallel use in mind and this RFC does not claim to address that broader
limitation.

\verysubsection{Failure mode and plugin loading order}

\noindent\colorbox{yellow!25}{\parbox{\dimexpr\linewidth-2\fboxsep}{%
  \textbf{Open Design Decision --- When does \texttt{H5Pset\_filter\_by\_name} load the plugin?}
  \texttt{H5Pset\_filter} does not load a plugin---it records the filter ID and
  \texttt{cd\_values} in the DCPL and defers all plugin loading to dataset creation
  or I/O time.  By analogy, \texttt{H5Pset\_filter\_by\_name} probably should not load
  the plugin at call time either.  If it does not, then name resolution,
  parameter parsing, and \texttt{set\_config} validation cannot occur at call time,
  and the failure mode described above (immediate \texttt{H5E\_NOFILTER} on a
  per-rank basis) is incorrect---failure would instead surface later, at
  \texttt{H5Dcreate} or I/O time, consistent with the existing plugin-loading model.

  The group must decide the plugin loading order for this API before the
  failure mode, error codes, and parallel guidance below can be specified
  normatively.  If plugin loading is deferred, the recommended pattern for
  pre-validating filter availability should be \texttt{H5Zfilter\_avail} (or the
  new \texttt{H5Zfilter\_avail\_by\_name}), which should be prominently documented
  as the correct way to check availability before committing to a DCPL
  configuration---something not currently recommended anywhere in the HDF5
  documentation.
}}

\subsection{VOL Connector Interaction}

Filters are applied to chunks by the native VOL connector during I/O. The new
string-based API resolves filter names and parameter strings into standard
\texttt{cd\_values} before modifying the DCPL. From the VOL connector's perspective,
the DCPL contains the same filter ID, flags, and \texttt{cd\_values} as if
\texttt{H5Pset\_filter} had been called directly.

Therefore, no changes to the VOL connector interface are required. Third-party
VOL connectors that intercept DCPLs will see standard filter configurations
regardless of whether they were set via the string API or the traditional API.

String-based filter configuration is resolved at \texttt{H5Pset\_filter\_by\_name}
call time, before any dataset is opened or created. It therefore has no
interaction with SWMR (which governs I/O access patterns on open datasets) or
with asynchronous VOL event sets (which are populated after the DCPL is fully
constructed).


%% ======================================================================
%%  Section 9 – CLI Tool Integration
%% ======================================================================
\section{CLI Tool Integration}
\label{sec:cli-tools}

\subsection{h5repack}

\texttt{h5repack} currently accepts filter specifications via the \texttt{-f} flag. The
existing syntax (\texttt{GZIP}, \texttt{SZIP}, \texttt{SHUF}, \texttt{FLET}, \texttt{NBIT},
\texttt{SOFF}, and \texttt{UD=id,flags,n,v1,...}) MUST continue to work unchanged.

A new form is added that accepts a filter name and parameter string:

\begin{Verbatim}[fontsize=\scriptsize]
    h5repack -f 'deflate=level=6' input.h5 output.h5
    h5repack -f 'zfp=mode=rate,rate=3.5' input.h5 output.h5
    h5repack -f '/group/dataset:zfp=mode=rate,rate=3.5' input.h5 output.h5
\end{Verbatim}

The equals sign is used as the delimiter between filter name and parameter
string. A colon cannot be used because it is already the separator between the
object path and filter specification in the existing \texttt{h5repack} syntax (e.g.,
\texttt{/group/dataset:GZIP=6}). Using equals for both the name-to-params delimiter
and within key=value pairs is unambiguous because the filter name is always
the text before the first equals sign.

The old filter name aliases used in \texttt{h5repack} (\texttt{GZIP}, \texttt{SHUF},
\texttt{FLET}, \texttt{SOFF}) are unrelated to the canonical names introduced in this
design. The old aliases MUST continue to work unchanged. Whether to also
register them as canonical name aliases (e.g., \texttt{"gzip"} $\rightarrow$
\texttt{H5Z\_FILTER\_DEFLATE}) is left as an implementation decision, but doing so
would simplify the \texttt{h5repack} parsing code by allowing it to delegate all
name resolution to \texttt{H5Pset\_filter\_by\_name}.

\verysubsection{\texttt{UD=} disambiguation}
\label{sec:ud-disambiguation}

The legacy user-defined filter syntax \texttt{UD=id,flags,n,v1,...} uses
\texttt{UD} as a reserved token.  To avoid ambiguity with the new
\textit{filtername}\texttt{=}\textit{params} syntax, \texttt{UD} is a
\emph{forbidden} canonical filter name and MUST be rejected by the name
registry.  The \texttt{h5repack} parser MUST always treat \texttt{UD=} as the
legacy syntax and handle disambiguation itself before calling
\texttt{H5Pset\_filter\_by\_name}---it is not the library's responsibility to
resolve this ambiguity.  Any filter that previously relied on the name
\texttt{UD} MUST use a different canonical name.

\subsection{h5dump}

The \texttt{PARAMS\_STRING} field is displayed only when \texttt{h5dump} is invoked
with \texttt{-p} (the existing flag that enables display of filter and property
details). Output without \texttt{-p} is unchanged. Because \texttt{PARAMS\_STRING} is
gated behind \texttt{-p}, and because format changes of this kind have been
accepted in prior major releases, no additional command-line flag is required.

\verysubsection{Normative output format (with \texttt{-p})}

When \texttt{-p} is active, for each filter in the pipeline that
implements \texttt{get\_config}, a \texttt{PARAMS\_STRING} line is appended immediately
after the existing filter line, indented to the same level as other filter
sub-fields:

\begin{verbatim}
    FILTERS {
      COMPRESSION DEFLATE { LEVEL 6 }
      PARAMS_STRING "level=6"
    }
\end{verbatim}

The value is the literal output of \texttt{get\_config} enclosed in double quotes.
If the string returned by \texttt{get\_config} contains a double-quote character,
it MUST be escaped as \texttt{\textbackslash\textquotedbl} in the h5dump output. For filters that do not implement \texttt{get\_config},
\texttt{H5Pget\_filter\_params\_by\_idx} provides a fallback string that formats the
raw \texttt{cd\_values} as a colon-separated list (see \SectionRef{sec:introspection}),
so \texttt{PARAMS\_STRING} can be displayed for all filters regardless of whether
they implement \texttt{get\_config}.

The format change is confined to the \texttt{FILTERS} block; all other
\texttt{h5dump} output is unaffected.

\end{document}
